\chapter{Normativa y legislación}

\section{Protección de datos}
El framework no almacena por sí mismo datos personales: es una herramienta de
construcción. No obstante, las aplicaciones que se construyan con él pueden
gestionar datos de carácter personal (identidades de usuario, autoría de
registros), por lo que el diseño se ha realizado teniendo presente el
Reglamento General de Protección de Datos (RGPD, Reglamento (UE) 2016/679) y
la Ley Orgánica 3/2018 de Protección de Datos Personales y Garantía de los
Derechos Digitales (LOPDGDD).

Las decisiones de diseño relevantes para este cumplimiento son: la
autenticación delega la verificación de credenciales en un proveedor externo
(Firebase) mediante verificación de \emph{token}, evitando almacenar
contraseñas; el modo de \emph{login} local existe únicamente para desarrollo y
no debe habilitarse en producción (RNF-07); el versionado con \gls{continuum}
mantiene un histórico de cambios que, en un despliegue real, debe acompañarse
de políticas de retención y del derecho de supresión; y el modelo de
autorización basado en roles y ACL permite implementar el principio de mínimo
privilegio sobre los datos.

\section{Licencias de software}
El proyecto \emph{UniWorks} se distribuye bajo licencia GNU General Public License (GPLv3), lo que permite su
reutilización en proyectos de la empresa sin fricción legal. Las dependencias
empleadas (FastAPI, SQLAlchemy, Pydantic, Alembic, SQLAlchemy-Continuum,
Angular, Formly, PostgreSQL, Redis, Docker) son software libre con licencias
permisivas (MIT, BSD, Apache~2.0, PostgreSQL License) compatibles entre sí y
con un uso comercial. No se ha incorporado código con licencias copyleft
fuertes que pudieran imponer restricciones a la distribución del framework.

\section{Propiedad intelectual}
El framework se desarrolla en colaboración con el \gls{itc}. La titularidad y
las condiciones de explotación se rigen por el acuerdo de colaboración entre
la universidad, la empresa y yo, y se reflejan en la
correspondiente solicitud de defensa del trabajo.
